% !TeX root = ../../Book1.tex
\section{紧算子及其对偶}
\label{b1:sec:C102}

\chapref{b1:ch:08}已给出紧算子的定义及其在复合、相加和算子范数极限下的稳定性。
本节补充两个后面需要的事实：无穷维空间的单位球不紧；紧性可以在算子与对偶算子之间传递。
前者用于把紧性转化成有限维性，后者使原方程与对偶方程能接受同样的分析。

\subsection{离开闭子空间的单位向量}
\label{b1:subsec:C10201}

\begin{lemma}[Riesz]\label{b1:lem:C1-riesz}
设 \(M\) 为赋范空间 \(X\) 的真闭线性子空间，\(0<\theta<1\)。
则存在单位向量 \(x\in X\)，使 \(\operatorname{dist}(x,M)>\theta\)。
\end{lemma}
\begin{proof}
取 \(z\notin M\)，并令 \(d=\operatorname{dist}(z,M)>0\)。
选取 \(m_0\in M\)，使 \(\|z-m_0\|<d/\theta\)，再置
\(x=(z-m_0)/\|z-m_0\|\)。对每个 \(m\in M\)，有
\[
 \|x-m\|
 =\frac{\bigl\|z-\bigl(m_0+\|z-m_0\|m\bigr)\bigr\|}{\|z-m_0\|}
 \geq\frac d{\|z-m_0\|}>\theta.
\]
取下确界即得所需估计。
\end{proof}

这就是\term[riesz-yinli]{Riesz 引理}[Riesz lemma]。
若 \(X\) 无穷维，可以依次对前面有限多个向量张成的子空间使用该引理，
构造单位向量列，使任意不同两项的距离大于 \(1/2\)。
有限维子空间闭，因此每一步都合法。该列无 Cauchy 子列，单位球便不可能相对紧。
结合有限维空间中有界列的子列定理，得到：赋范空间的闭单位球紧，当且仅当空间有限维。
这里没有假定空间存在内积，也没有使用正交向量。

这个引理中的 \(\theta<1\) 留出了不一定取到最短距离的余量。在 Hilbert 空间中
可以取与 \(M\) 正交的单位向量而得到距离恰为 \(1\)；一般 Banach 空间不能靠“取正交补”
完成这一步。

\subsection{紧性在对偶中的传递}
\label{b1:subsec:C10202}

\begin{theorem}[Schauder]\label{b1:thm:C1-schauder}
设 \(X,Y\) 为 Banach 空间，\(T\in\mathcal B(X,Y)\)。
则 \(T\) 紧，当且仅当 \(T':Y^*\rightarrow X^*\) 紧。
\end{theorem}
\begin{proof}
先设 \(T\) 紧。令 \(C=\overline{T(B_X(0,1))}\)，这是 \(Y\) 中的紧集。
考虑 \(Y^*\) 的闭单位球中一列泛函 \(g_n\)。对每个正整数 \(k\)，选取 \(C\) 的有限
\(1/k\)-网。所有这些网点组成至多可数的集合 \(\{y_j\}\)。
由于 \(|g_n(y_j)|\leq\|y_j\|\)，逐点抽子列并作对角选取，可使
\(g_{n_\ell}(y_j)\) 对每个 \(j\) 收敛。

给定 \(\varepsilon>0\)，选一个网的半径小于 \(\varepsilon/4\)。
在有限个网点上同时使用 Cauchy 条件，便有：对充分大的 \(\ell,m\)，
\[
 \sup_{y\in C}|g_{n_\ell}(y)-g_{n_m}(y)|<\varepsilon.
\]
事实上，先将 \(y\) 移到某个网点，两个泛函各造成不超过网半径的误差；
网点处的差小于 \(\varepsilon/2\)。于是
\[
 \|T'g_{n_\ell}-T'g_{n_m}\|
 =\sup_{\|x\|\leq1}|g_{n_\ell}(Tx)-g_{n_m}(Tx)|<\varepsilon.
\]
连续对偶 \(X^*\) 完备，故这条像子列收敛。任意有界泛函列可先缩放到单位球，
所以 \(T'\) 紧。

反过来，若 \(T'\) 紧，将已证方向用于 \(T'\)，得到 \(T''\) 紧。
设 \((x_n)\) 是 \(X\) 中的有界列。等距性使 \((J_Xx_n)\) 有界，因而可取子列，
使
\[
 T''J_Xx_{n_k}=J_YTx_{n_k}
\]
在 \(Y^{**}\) 中收敛。等距嵌入 \(J_Y(Y)\) 由于 \(Y\) 完备而闭：
其中 Cauchy 列的原像在 \(Y\) 中也为 Cauchy 列，极限仍属于该像。
所以所得极限属于 \(J_Y(Y)\)，原像列 \(Tx_{n_k}\) 在 \(Y\) 中收敛，证明 \(T\) 紧。
\end{proof}

\term[schauder-jinxing]{Schauder 定理}[Schauder theorem]不要求 \(X,Y\) 自反。
反向证明虽然进入了二次对偶，但只使用自然嵌入的等距性和闭像，不能把这一步改写成
“因为 \(Y=Y^{**}\)”。

\subsection{Hilbert 空间中的有限秩逼近}
\label{b1:subsec:C10203}

\chapref{b1:ch:08}证明过用有限秩算子作算子范数逼近足以推出紧性。
在一般 Banach 空间中，逆向结论需要额外条件；若值域是 Hilbert 空间，则有如下直接构造。

\begin{proposition}\label{b1:prop:C1-hilbert-finite-rank}
设 \(X\) 为 Banach 空间，\(H\) 为 Hilbert 空间，\(T:X\rightarrow H\) 紧。
给定 \(\varepsilon>0\)，存在有限维子空间 \(E\subset H\)，使其正交投影 \(P_E\) 满足
\[
 \|T-P_ET\|<\varepsilon.
\]
\end{proposition}
\begin{proof}
在紧集 \(\overline{T(B_X(0,1))}\) 中取有限个半径小于 \(\varepsilon/2\) 的球覆盖，
球心为 \(y_1,\ldots,y_m\)，令 \(E\) 为这些向量张成的空间。
对 \(\|x\|\leq1\)，可取一个球心使 \(\|Tx-y_j\|<\varepsilon/2\)。
正交投影给出最近点，故
\[
 \|Tx-P_ETx\|=\operatorname{dist}(Tx,E)
 \leq\|Tx-y_j\|<\varepsilon/2.
\]
取单位球上的上确界便得到
\(\|T-P_ET\|\leq\varepsilon/2<\varepsilon\)。
\end{proof}

有限维近似由紧像的有限网确定，而不是由 \(X\) 的某个预先给定的坐标截断确定。
即使 \(H\) 不可分，每次所需的近似空间也只有有限维。又因为
\((P_ET)'=T'P_E'\)，这个构造与 Schauder 定理相容，但它没有证明任意紧算子都自伴，
更没有把有限秩近似误当成特征向量分解。
